For clarity, the inherited scalar argument still applies to the new
indices. At each odd contact prime of $\beta$ outside $S_*$, the
constant-left scalar norm tree excludes negative odd raw valuation and
gives all six assigned coefficient orders. The first four put $S,J_0$
in the formal group at every place above that prime. Writing
$2S=sP_*$ and $2J_0=j_0P_*$, with nonzero integer $j_0$, the last two
give $v_p(s/j_0)\ge32$, hence $v_p(s)\ge32$. The logarithm comparison
uses the exact identity $[2j_0]S=[2s]J_0$. It is valid also at inert
primes: both rational coefficients are tested and the extension is
unramified outside two. The unchanged scalar raw polynomial and its
small-prime correction detect every prime of $M_5(\beta)$, including
an odd infinite contact. Consequently
\[
 M_5(\beta)^{32}\mid s,\qquad
 h(\beta)\le H M_5(\beta)^{19}\le H|s|^{19/32}.
\]
The height theorem's domain condition holds. The value $\beta=X^2+2$
is none of $0,1,-343/32$, and $\beta=2401$ would require the nonsquare
integer $2399$ to be a rational square. Doubling both indices therefore
changes neither the scalar ratio nor its divisibility argument.
